Study~4 registered each family with the human finding it would be compared with.
Table~\ref{tab:human} sets the results beside those findings. The comparison is of direction only:
the agents face no real stakes, every agent is one model, and none of the quantities is calibrated
to a human measure.

\begin{table}[t]
\centering
\small
\caption{Results set against the human findings they were registered to be compared with.}
\label{tab:human}
\begin{tabularx}{\linewidth}{@{}L L L >{\raggedright\arraybackslash}p{2.1cm}@{}}
\toprule
Human finding & Our analogue & Result & Same direction? \\
\midrule
Identifying individuals raises cooperation \citep{andreoni2004public,christens2019identification} &
Any stable individual handle vs none: pair hunts (A1, A2; Studies 1--2) &
Pair hunting rose two- to fivefold with any handle & Yes \\
\addlinespace
Identifying only groups does little \citep{christens2019identification} &
A band without an individual handle (D5, E2) &
No change in pair hunting or the commons & Yes \\
\addlinespace
Names alone have no direct effect on public-goods contributions \citep{sampaio2023identity} &
Named vs stable ID: the commons (E3) &
No difference, at a floor & Consistent, not informative \\
\addlinespace
Arbitrary labels produce in-group favouritism \citep{tajfel1971social} &
In-band couple, hunt and gift bias (D1--D3); endowment to an other-band child (exploratory) &
In-band couples and hunts; less given to an other-band child & Yes \\
\addlinespace
Shared group labels split LLM agents into factions and degrade coordination \citep{delgiudice2026identity} &
A band added to an individual ID (D4) &
Pair hunting doubled, mostly within bands & No, opposite \\
\addlinespace
Children named after a parent receive more of that parent's investment \citep{saffa2021namesake} &
Namesake parent minus co-parent (C1--C3) &
No difference at birth; weak lead in later gifts & No \\
\addlinespace
First-round public-goods contributions of 40--60\%, declining with repetition \citep{ledyard1995public,chaudhuri2011sustaining} &
Commons contribution share by generation &
About 3\%, rising from 1\% to 5\% & No \\
\addlinespace
Help concentrates on kin \citep{hamilton1964genetical} &
Share of gift amounts to kin (Studies 1 and 4, exploratory) &
Most concentrated without names (84--90\%) & Yes \\
\bottomrule
\end{tabularx}
\end{table}

Four of the eight comparisons agree in direction, one is uninformative, and three disagree. The
agreements cluster around identification and labels: individual handles enable cooperation, group
labels alone do not, and arbitrary group labels bias choices toward the in-group. The
disagreements concern investment and contribution: namesaking does not move parental investment,
and the agents free-ride on the commons almost completely. The commons result carries a caveat
specific to this world. Survival is decided by rank, so anything put into the pot benefits rivals
as much as the contributor, and contributing always lowers an agent's standing relative to those
who hold back. A standard public-goods game has no such tournament.
